% !TeX root=main.tex
\section{تکسونومی روش‌های مقیاس‌پذیری}

برای سامان‌دهی پیشینه پژوهش، روش‌های مقیاس‌پذیری را در چهار خانواده اصلی دسته‌بندی می‌کنیم. شکل \ref{fig:scaling-taxonomy} نمای کلی این تکسونومی را نشان می‌دهد.

\begin{figure}[h]
\centering
\begin{tikzpicture}[
  level 1/.style={sibling distance=5.2cm, level distance=1.6cm},
  level 2/.style={sibling distance=2.4cm, level distance=1.4cm},
  every node/.style={rectangle, draw, rounded corners, align=center, font=\small, inner sep=4pt}
]
\node {مقیاس‌پذیری بلاکچین}
  child {node {لایه اول\\$\text{(L1)}$}}
    child {node {افزایش\\ظرفیت بلوک}}
    child {node {بلوک\\سریع‌تر}}
    child {node {تغییر\\اجماع $\text{(PoS)}$}}
  }
  child {node {لایه دوم\\$\text{(L2)}$}}
    child {node {کانال\\حالت}}
    child {node {رول‌آپ}}
  }
  child {node {موازی‌سازی}
    child {node {شاردینگ}}
    child {node {چندزنجیره\\/ زنجیره جانبی}}
  };
\end{tikzpicture}
\caption{تکسونومی روش‌های مقیاس‌پذیری: هر شاخه نشان‌دهنده «دیدگاه طراحی» متفاوت است}
\label{fig:scaling-taxonomy}
\end{figure}

\begin{table}[h]
\centering
\caption{تکسونومی روش‌های مقیاس‌پذیری و مکانیزم افزایش \lr{TPS}}
\label{tab:scaling-taxonomy}
\small
\begin{tabular}{|p{2.4cm}|p{3.2cm}|p{3.8cm}|p{2.8cm}|}
\hline
\textbf{خانواده} & \textbf{نمونه‌های شاخص} & \textbf{مکانیزم مقیاس‌پذیری} & \textbf{دیدگاه طراحی} \\
\hline
\lr{L1} -- ظرفیت بلوک & \lr{BCH}, \lr{BSV}, \lr{SegWit} & افزایش تعداد تراکنش در هر بلوک & سادگی؛ هزینه نود بالا \\
\hline
\lr{L1} -- زمان بلوک & برخی \lr{altcoins} & کاهش فاصله بلوک & تأخیر کمتر؛ ریسک امنیت \\
\hline
\lr{L1} -- اجماع & \lr{Ethereum PoS}, \lr{Cardano} & جایگزینی \lr{PoW} با \lr{PoS}/\lr{DPoS} & کارایی انرژی؛ ریسک تمرکز سهام \\
\hline
\lr{L2} -- کانال & \lr{Lightning}, \lr{Raiden}, \lr{Plasma} & تراکنش خارج از زنجیره + تسویه دوره‌ای & \lr{off-chain}؛ وابستگی نقدینگی \\
\hline
\lr{L2} -- رول‌آپ & \lr{Optimistic}/\lr{ZK Rollup} & اثبات دسته‌ای روی \lr{L1} & امنیت \lr{L1}؛ پیچیدگی اثبات \\
\hline
موازی‌سازی -- شارد & \lr{Ethereum 2.0}, \lr{Zilliqa} & تقسیم حالت/تراکنش بین شاردها & \lr{TPS} بسیار بالا؛ پیچیدگی بین‌شاردی \\
\hline
موازی‌سازی -- چندزنجیره & \lr{RSK}, \lr{Liquid}, \textbf{معماری پیشنهادی} & زنجیره‌های مستقل + لنگر در \lr{L1} & تخصصی‌سازی + پردازش موازی \\
\hline
\end{tabular}
\end{table}

\section{زنجیره‌های جانبی}

زنجیره‌های جانبی \lr{(Sidechains)} بلاکچین‌هایی هستند که به موازات بلاکچین اصلی عمل می‌کنند و امکان انتقال دارایی بین زنجیره اصلی و زنجیره جانبی را فراهم می‌کنند.

\subsection{مفهوم زنجیره جانبی}

زنجیره جانبی یک بلاکچین مستقل است که از طریق یک پروتکل دوطرفه به بلاکچین اصلی متصل می‌شود. این اتصال امکان انتقال دارایی بین دو زنجیره را فراهم می‌کند.

\subsection{نمونه‌های عملی}

\subsubsection{RSK \lr{(Rootstock)}}

RSK یک زنجیره جانبی برای بیتکوین است که امکان اجرای قراردادهای هوشمند را فراهم می‌کند \cite{rsk2015}. RSK از merge mining استفاده می‌کند که به ماینرها اجازه می‌دهد همزمان بیتکوین و RSK را استخراج کنند.

\subsubsection[\protect\hbox{\lr{Liquid Network}}]{\texorpdfstring{\lr{Liquid Network}}{Liquid Network}}

\lr{Liquid Network} یک زنجیره جانبی برای بیتکوین است که توسط \lr{Blockstream} توسعه یافته است. این شبکه امکان تراکنش‌های سریع و محرمانه را فراهم می‌کند.

زنجیره‌های جانبی انعطاف‌پذیری در طراحی پروتکل، امکان آزمایش ویژگی‌های جدید بدون تغییر در زنجیره اصلی و تخصصی‌سازی برای کاربردهای خاص را فراهم می‌کنند؛ با این حال نیاز به پروتکل‌های پیچیده برای ارتباط با زنجیره اصلی، ریسک امنیتی در صورت نقص در پروتکل ارتباطی و در برخی پیاده‌سازی‌ها نیاز به اعتماد به نودهای فدرال \lr{(federated nodes)} دارند.

\section{الگوریتم‌های اجماع جایگزین}

\subsection{اثبات سهام \lr{(Proof of Stake)}}

اثبات سهام یک الگوریتم اجماع است که در آن اعتبارسنج‌ها بر اساس میزان سهامی که دارند انتخاب می‌شوند. این الگوریتم انرژی کمتری نسبت به PoW مصرف می‌کند.

اثبات سهام مصرف انرژی بسیار کمتر، سرعت بالاتر در تأیید و مقیاس‌پذیری بهتر دارد؛ اما ریسک تمرکز در صورت توزیع نابرابر سهام، مشکل Nothing at Stake و نیاز به مکانیزم‌های پیچیده برای امنیت از معایب آن است. پروژه‌هایی مانند Decred و Horizen از ترکیب PoW و PoS استفاده می‌کنند.

\section{راه‌حل‌های افزایش اندازه بلوک}

\subsection[\protect\hbox{\lr{Bitcoin Cash}}]{\texorpdfstring{\lr{Bitcoin Cash}}{Bitcoin Cash}}

\lr{Bitcoin Cash} در سال \lr{2017} با افزایش اندازه بلوک از \lr{1} مگابایت به \lr{8} مگابایت (و بعداً \lr{32} مگابایت) ایجاد شد. این تغییر باعث افزایش \lr{TPS} می‌شود اما نیاز به منابع بیشتری دارد.

افزایش اندازه بلوک TPS را به صورت ساده و با تغییرات کم در پروتکل بالا می‌برد؛ اما نیاز به منابع بیشتر برای ذخیره‌سازی و پردازش، افزایش زمان همگام‌سازی برای نودهای جدید و ریسک تمرکز در صورت افزایش بیش از حد اندازه بلوک از معایب آن است.

\section{مقایسه راه‌حل‌ها}

جدول \ref{tab:solutions_comparison} مقایسه کیفی راه‌حل‌های مرورشده را نشان می‌دهد. ستون «رشد مقیاس» بیان می‌کند آیا با افزودن منابع/زیرسیستم، \lr{TPS} به‌صورت تقریباً خطی رشد می‌کند یا خیر.

\begin{table}[h]
\centering
\caption{مقایسه راه‌حل‌های مقیاس‌پذیری (کیفی)}
\label{tab:solutions_comparison}
\small
\begin{tabular}{|p{2.6cm}|p{1.6cm}|p{1.5cm}|p{1.5cm}|p{1.8cm}|p{2.2cm}|}
\hline
\textbf{راه‌حل} & \textbf{TPS تقریبی} & \textbf{امنیت} & \textbf{غیرمتمرکز} & \textbf{رشد مقیاس} & \textbf{نیاز به تغییر \lr{L1}} \\
\hline
بیتکوین (پایه) & $\sim$7 & $\star\star\star\star\star$ & $\star\star\star\star\star$ & ثابت & --- \\
\hline
افزایش اندازه بلوک & 50--100 & $\star\star\star\star$ & $\star\star\star\star$ & محدود & کم \\
\hline
شبکه لایتنینگ & $10^3+$ & $\star\star\star$ & $\star\star\star$ & وابسته به کانال & خیر (\lr{L2}) \\
\hline
رول‌آپ (\lr{L2}) & $10^3$--$10^4$ & $\star\star\star\star$ & $\star\star\star\star$ & با ظرفیت \lr{prover} & قرارداد/پروتکل \\
\hline
شاردینگ (اتریوم ۲) & $10^4$+ & $\star\star\star\star$ & $\star\star\star\star$ & خطی با شارد & زیاد \\
\hline
زنجیره جانبی (\lr{RSK}/\lr{Liquid}) & متغیر & $\star\star\star\star$ & $\star\star\star$--$\star\star\star\star$ & محدود (تعداد ثابت) & پل دوطرفه \\
\hline
\textbf{معماری پیشنهادی} & خطی با $N$ & $\star\star\star\star\star$ (\lr{L1}) & $\star\star\star\star$ & \textbf{خطی با $N$} & کم (لایه جانبی) \\
\hline
\end{tabular}
\end{table}

\subsection{جایگاه معماری پیشنهادی}

برخلاف لایتنینگ که مقیاس را به «کانال‌های خصوصی» واگذار می‌کند، و برخلاف شاردینگ که نیازمند بازطراحی کامل پروتکل است، معماری این پایان‌نامه \textbf{پردازش موازی روی چند زنجیره \lr{PoW} سبک‌وزن} را با \textbf{لنگر امنیتی \lr{PoW} در زنجیره اصلی} ترکیب می‌کند. از این نظر به خانواده «موازی‌سازی -- چندزنجیره» نزدیک است، اما تمرکز آن بر \lr{TPS} قابل پیش‌بینی و سازگاری با مدل \lr{UTXO} بیتکوین است (فصل سوم).

\section{خلاصه و نتیجه‌گیری}

در این فصل، مسئله مقیاس‌پذیری را به‌عنوان تضاد میان ظرفیت بلوک، تأخیر اجماع و هزینه نود تحلیل کردیم. تکسونومی شکل \ref{fig:scaling-taxonomy} و جدول \ref{tab:scaling-taxonomy} نشان داد که هر روش از یک «دیدگاه طراحی» متفاوت پیروی می‌کند: افزایش ظرفیت \lr{L1}، اجرای \lr{off-chain}، شاردینگ، یا چندزنجیره‌ای. راه‌حل‌های \lr{L2} \lr{TPS} بالایی ارائه می‌دهند اما پیچیدگی عملیاتی و وابستگی نقدینگی دارند؛ شاردینگ ظرفیت بسیار بالایی هدف می‌گیرد اما پیاده‌سازی و امنیت بین‌شاردی دشوار است؛ زنجیره‌های جانبی کلاسیک اغلب تعداد محدودی زنجیره دارند. معماری پیشنهادی با \textbf{افزودن خطی زنجیره‌های جانبی} و \textbf{حفظ \lr{PoW} در زنجیره اصلی}، شکاف «مقیاس خطی + سازگاری با بیتکوین» را هدف می‌گیرد که در فصل سوم به‌صورت رسمی تعریف می‌شود.
